\chapter{Modellek}
\label{ch:models}

      % pl. Peano aritmetika (ahogy fent van megadva)
      %   relar 2 = 𝟚, false az egyenloseg, Rel {Γ}{2} false : Tm Γ × Tm Γ × ⊤ → For Γ
      %   zero := fun {Γ}{0} tt tt : Tm Γ
      %   suc (n : Tm Γ) := fun {Γ}{1} tt (n, tt) : Tm Γ
      %   Eq (u v : Tm Γ) := Rel {n = 2} false (u , v , tt)
      % (x,p:x=3) ⊢ ∀y.y=3 ⊃ y=x
      % ? : Pf (◇ ▹ₜ ▹ₚ Eq qₜ (suc (suc (suc zero)))) (Forall (Eq qₜ (suc (suc (suc zero))) ⊃ Eq qₜ (qₜ[pₜ][pₚ])))


A szemantikai modellek a megadott elsőrendű logika jelentését, értelmezését adják meg. A három vizsgált modell - a Bool modell, a Tarski modell és a Family modell - különböző megközelítéseket kínál a logikai kifejezések értelmezésére és kiértékelésére.

\begin{itemize}
\item Bool modell: A legegyszerűbb és legközvetlenebb megközelítés, amely a logikai kifejezéseket az igaz és hamis értékekre képezi le, valamint a logikai műveletek a Bool értékek közötti műveletekkel azonosak. Ez a klasszikus kétértékű szemantika.

\item Tarski modell: A logikai kifejezések értelmezését a halmazelmélet alapjaira helyezi. Ebben a modellben a logikai változók és kifejezések értékei halmazok, illetve halmazok közötti relációk. A Tarski szemantika a logikai igazságot a kifejezések halmazelméleti interpretációjával határozza meg, amely lehetővé teszi a logikai állítások formális és precíz vizsgálatát.

\item Family modell: Egy indexelt családok segítségével felépített szemantika. Felfogható a Kripke-szemantika egy specializációjaként. Ez a modell, a Kripke modellel szemben, nem specifikál explicit elérhetőségi relációt az indexek között, ami egyszerűsíti a struktúrát.
\end{itemize}


\section{Bool modell}


A Bool modellhez közvetlenül az Agda beépített Bool típusát használjuk, így a logikai összekötők (∧, ∨, →, ¬) is a megfelelő Bool műveletekre fordítódnak. A formulák kiértékelése során minden atomi állítás, predikátum és függvény végső soron egy Bool értékre redukálódik. Az univerzális és egzisztenciális kvantorokra gondolhatunk úgy, mint listák feletti all és any műveletek.


\section{Tarski modell}


A Tarski modell az elsőrendű logika standard szemantikai modellje. Egy nem üres domain (univerzum) felett értelmezett interpretáció.

Az Agdában implementált Tarski modell a következő komponensekből áll:

Egy tetszőleges típus, amely a domaint reprezentálja
Függvények, amelyek a logikai nyelv szimbólumait a domain elemeire, függvényeire és relációira képezik le
Változókiértékelő függvény, amely a változókat a domain elemeire képezi le
Rekurzív interpretációs függvény, amely a formulákat a domain feletti állításokra képezi le



\section{Family modell}

A Family modell felfogható a Kripke-szemantika egy specializációjaként.

Indexelés mint lehetséges világok: A Family modellben szereplő I indexhalmaz elemei megfeleltethetők a Kripke szemantika lehetséges világainak.
Az elérhetőségi reláció implicit módon adott vagy triviális
A D(i) tartományok explicit módon specifikáltak minden i indexre
A logikai operátorok szemantikája hasonlóan van definiálva, de a modell szerkezete egyszerűbb és célorientáltabb

Ahogyan a Kripke-modellekben minden lehetséges világban kiértékeljük a formulákat, a Family modellben minden indexhez tartozó struktúrán értelmezzük a kifejezéseket és formulákat.
Elérhetőségi reláció hiánya: A Kripke-modellekben alapvető szerepet játszik a világok közötti elérhetőségi reláció (R), amely meghatározza a modális operátorok szemantikáját. A Family modell nem specifikál explicit elérhetőségi relációt az indexek között, ami egyszerűsíti a struktúrát.


Minden predikátum egy típuscsalád, amely a domain elemeihez típusokat rendel
A logikai konnektívák típuskonstruktorokként implementálódnak (pl. függő összeg és szorzat típusok)
A kvantorok függő típusokként implementálódnak (Π és Σ típusok)
A bizonyítások a megfelelő típusok értékei/lakói



\section{Modellek összehasonlítása}

A modellek összehasonlítása

Expresszivitás

Bool modell: A legkevésbé expresszív, csak véges domainek felett használható hatékonyan, és a klasszikus kétértékű logikára korlátozódik.
Tarski modell: Nagyon expresszív, a klasszikus elsőrendű logika teljes szemantikáját képes modellezni, beleértve a végtelen domaineket is.
Family modell: Expresszív és konstruktív jellegű, de más filozófiai alapokon nyugszik, mint a klasszikus logika.

Implementációs komplexitás

Bool modell: Egyszerű implementáció, könnyen érthető kód.
Tarski modell: Közepes komplexitás, több absztrakciós réteget igényel.
Family modell: A legösszetettebb implementáció, amely mély típuselméleti ismereteket igényel.

Alkalmazási területek

Bool modell: Ideális tautológiák ellenőrzésére, véges modell ellenőrzésre és egyszerű példákra.
Tarski modell: Alkalmas általános célú logikai modellezésre, tételbizonyítókhoz és modellellenőrzőkhöz.
Family modell: Kiválóan alkalmas bizonyítássegédekhez, programlevezetéshez és verifikációhoz.


\section{Összegzés}

A három bemutatott modell különböző megközelítéseket kínál az elsőrendű logika szemantikájának implementálására Agdában. Mindegyik modellnek megvannak a maga erősségei és gyengeségei, valamint jellemző alkalmazási területei:

A Bool modell egyszerűsége és hatékonysága miatt alkalmas gyors prototípusok készítésére és véges domainek feletti ellenőrzésekre.
A Tarski modell matematikai precizitása és általánossága miatt a legjobb választás a klasszikus elsőrendű logika teljes szemantikájának implementálására.
A Family modell típuselméleti megalapozottsága miatt különösen alkalmas olyan alkalmazásokhoz, ahol a bizonyítások és számítások integrációja fontos.

Az implementációnkban mindhárom modellt megvalósítottuk, demonstrálva ezzel az Agda típusrendszerének erejét és rugalmasságát. A modellek közötti formális kapcsolatokat is definiáltuk, bizonyítva, hogy bizonyos feltételek mellett konzisztens eredményeket adnak.